\section{Conclusion and Future Work}
\label{sec:conclusion}

This paper presented CardioSentinel, a privacy-preserving, one-shot federated model for the multi-label
classification of 12-lead ECGs into the five PTB-XL diagnostic superclasses. Four clients trained a
1D CNN, a 1D ResNet, a BiLSTM and a CNN-LSTM on disjoint, patient-separated quarters of the training
data, calibrated them, and shared only their class probabilities with a server that combines them on
the probability simplex. A parsimony-first selection on patient-grouped validation predictions chose
equal weighting over learned static weights and a per-patient dynamic gate in all three seeds. On the
untouched test fold the federated model reached a macro AUROC of 0.905 $\pm$ 0.000, a macro F1 of
70.7\% and a label-wise accuracy of 86.9\%, improved on the best single client in every seed
(+0.009 to +0.012 AUROC, paired bootstrap), and came within 0.004 AUROC of a model trained on all
data pooled together while matching its macro F1. It outperformed every model trainable at a single
site, including hand-crafted-feature XGBoost, whereas FedAvg with a shared CNN-LSTM was 0.002 AUROC
better after 20 rounds and pooled-data XGBoost 0.004 better; the proposed scheme therefore trades a small
loss in accuracy for a single communication round and a free choice of architecture. Client-subset ablations and Shapley values showed
that every institution contributes, which is why equal weighting outperformed learned weights that
discard the weaker clients. Integrated Gradients attributed the predictions to leads and waveform
segments, and per-finding triage bands automated 53.7\% of the findings at a precision of 0.937 and an
accuracy of 0.967 while referring the remainder to a cardiologist.

Future work will follow four directions. First, the method will be evaluated on genuinely
heterogeneous institutions with different devices, populations and label distributions, where the
diversity of the clients and the need for a learned gate are expected to be larger. Second, a client
trained on hand-crafted amplitude features, such as the XGBoost model that reached 0.909 AUROC on pooled
data, will be added to the federation, and the parameter-averaging baselines will be repeated over
several seeds and with longer training. Third, the gate will receive inputs that actually differ between
clients, such as additional modalities or site metadata, so that per-patient weighting can become
informative, and formal guarantees such as differential privacy on the shared probabilities will be
studied. Fourth, the recognition of hypertrophy will be improved, for example by restoring absolute
voltage information that per-recording normalisation removes, and the explanations will be validated
with cardiologists.
